%%%%%%%%%%%%%%%%%%%%%%%%%%%%%%%%%%%%%%%%%%%%%%%%%%%%%%%%%%%%%%%%%%%%%%%%%%%%
%% MCM/ICM paper starter — built on the official COMAP MCM-ICM_Summary.tex.
%% This skeleton already encodes the layout lessons the Kit enforces:
%%   * official 3-column Summary Sheet header (Problem / Summary Sheet / Team)
%%   * title-first Summary (title, then a centered "Summary" heading, then body)
%%   * hyperref with hidelinks (no colored link boxes)
%%   * captions BELOW floats, small font, bold label, width-limited
%%   * float parameters tuned so figures land on text pages, not sparse float pages
%%   * an AI Use Report AFTER references (excluded from the 25-page count)
%% Replace the bracketed guidance with real content. Aim to FILL ~20-25 pages
%% of solution with substance the problem rewards (never filler).
%%%%%%%%%%%%%%%%%%%%%%%%%%%%%%%%%%%%%%%%%%%%%%%%%%%%%%%%%%%%%%%%%%%%%%%%%%%%
\documentclass[12pt]{article}
\usepackage{geometry}
\geometry{left=1in,right=0.75in,top=1in,bottom=1in}

%%%%%%%%%%%%%%%%%%%%%%%%%%%%%%%%%%%%%%%%
% Replace ABCDEF with your chosen problem and 1111111 with your Team Control Number.
\newcommand{\Problem}{ABCDEF}
\newcommand{\Team}{1111111}
%%%%%%%%%%%%%%%%%%%%%%%%%%%%%%%%%%%%%%%%

\usepackage{newtxtext}
\usepackage{amsmath,amsthm}   % load amsmath BEFORE newtxmath; do NOT load amssymb (\Bbbk clash)
\usepackage{newtxmath}        % must come after amsmath
\usepackage[pdftex]{graphicx}
\usepackage{xcolor}
\usepackage{booktabs}
\usepackage{fancyhdr}
\usepackage{caption}
\usepackage[hidelinks]{hyperref}   % hidelinks => no colored link boxes in the PDF

% Captions: below floats, small font, bold label, width-limited (not full text width).
\captionsetup{font=small,labelfont=bf,width=0.86\textwidth}

% Float parameters: force floats onto text pages and prevent near-empty float-only pages.
\renewcommand{\topfraction}{0.92}
\renewcommand{\bottomfraction}{0.85}
\renewcommand{\textfraction}{0.06}
\renewcommand{\floatpagefraction}{0.88}
\setcounter{totalnumber}{4}

\lhead{Team \Team}
\rhead{}
\cfoot{}

\newtheorem{theorem}{Theorem}
\newtheorem{definition}{Definition}

\begin{document}
\graphicspath{{../figures/generated/}{../figures/concept/}{./}}
\DeclareGraphicsExtensions{.pdf,.png,.jpg}

%%%%%%%%%%% Summary Sheet header (official) %%%%%%%%%%%
\thispagestyle{empty}
\vspace*{-16ex}
\centerline{\begin{tabular}{*3{c}}
    \parbox[t]{0.3\linewidth}{\begin{center}\textbf{Problem Chosen}\\ \Large \textcolor{red}{\Problem}\end{center}}
    & \parbox[t]{0.3\linewidth}{\begin{center}\textbf{2026\\ MCM/ICM\\ Summary Sheet}\end{center}}
    & \parbox[t]{0.3\linewidth}{\begin{center}\textbf{Team Control Number}\\ \Large \textcolor{red}{\Team}\end{center}} \\
    \hline
\end{tabular}}

%%%%%%%%%%% Title-first Summary (title, THEN "Summary" heading, THEN body) %%%%%%%%%%%
\begin{center}
    {\Large\bfseries [Paper Title: a specific, results-forward title]}\\[2ex]
    \textbf{Summary}
\end{center}

% One page. Topic sentences and numbers, not generic prose. A judge should grasp the
% value of the solution from this page alone: problem -> model family & why -> data/params
% -> headline quantitative results -> decisions/recommendations -> key limitations.
[Summary body. Lead with the problem in one sentence. Name the model family and why it
fits. State the data/parameters. Give the headline numbers. State the recommendation.
Close with the main limitation. Fill the page.]

\vspace{1ex}
\noindent\textbf{Keywords:} [keyword; keyword; keyword]

%%%%%%%%%%%%%%%%%%%%%%%%%%%%%%
\clearpage
\pagestyle{fancy}
\tableofcontents
\newpage
\setcounter{page}{1}
\rhead{Page \thepage\ }
%%%%%%%%%%%%%%%%%%%%%%%%%%%%%%

\section{Introduction}
\subsection{Problem Background and Restatement}
[Restate the problem tied to the specific tasks the contest asks for.]

\subsection{Our Approach}
[One paragraph overview + the concept/workflow figure below.]

% Concept figure: register it in figures/concept_src/<name>.json and add the file to
% diagram_sources in workflow_config.json so the diagram-quality gate inspects it.
% \begin{figure}[htbp]
%   \centering
%   \includegraphics[width=0.86\textwidth]{model_flow.png}
%   \caption{Solution workflow: data inputs, parameter estimation, model, solver,
%            validation, and outputs, with the feedback loop used to calibrate.}
%   \label{fig:workflow}
% \end{figure}

\section{Assumptions and Justifications}
[Assumptions with consequences, not boilerplate. Each: assumption -> why -> effect.]

\section{Notation}
[A notation table only if it saves the reader time.]

\section{Model Development}
\subsection{Data and Parameter Provenance}
[Where each number comes from: cited source or data-derived. Be honest.]

\subsection{Derivation}
[Governing equations with every symbol defined; the solution method; algorithm.]

\section{Results}
[Headline numbers prominent and traceable to key\_results.csv. Figures/tables carry
the argument. Each figure earns its place.]

\section{Validation and Sensitivity}
[Validation or sanity checks against data or a baseline; sensitivity + uncertainty with
consequences. Report honest residuals.]

\section{Discussion}
\subsection{Strengths and Weaknesses}
\subsection{Recommendations}

\section{Conclusion}
[Tie results back to the tasks.]

% References auto-generate the section; make it show in the ToC.
\addcontentsline{toc}{section}{References}
\begin{thebibliography}{99}
\bibitem{ref1} [Author, \emph{Title}, Source, Year.]
\end{thebibliography}

% ---- AI Use Report: AFTER references; does NOT count toward the 25-page limit. ----
\clearpage
\addcontentsline{toc}{section}{Report on the Use of AI}
\section*{Report on the Use of AI}
[Disclose AI tools and how they were used, per the current COMAP rules.]

\end{document}
